\subsection{Dividing FPSs}

\subsubsection{Conventions}

We shall make ourselves at home in the ring $K\left[  \left[  x\right]
\right]  $ a bit more. (Recall that $K$ is a fixed commutative ring.)

\begin{convention}
\label{conv.fps.const}From now on, we identify each $a\in K$ with the constant
FPS $\underline{a}=\left(  a,0,0,0,0,\ldots\right)  \in K\left[  \left[
x\right]  \right]  $.
\end{convention}

This is motivated by the fact that $\underline{a}=a+0x+0x^{2}+0x^{3}+\cdots$
for any $a\in K$. Convention \ref{conv.fps.const} does not cause any dangerous
ambiguities, because we have%
\begin{align*}
\underline{a+b}  &  =\underline{a}+\underline{b}\ \ \ \ \ \ \ \ \ \ \text{and}%
\\
\underline{a-b}  &  =\underline{a}-\underline{b}\ \ \ \ \ \ \ \ \ \ \text{and}%
\\
\underline{a\cdot b}  &  =\underline{a}\cdot\underline{b}%
\ \ \ \ \ \ \ \ \ \ \text{for any }a,b\in K
\end{align*}
(check this!), and because the zero and the unity of the ring $K\left[
\left[  x\right]  \right]  $ are $\underline{0}$ and $\underline{1}$, respectively.

Furthermore, I will stop using boldfaced letters (like $\mathbf{a}%
,\mathbf{b},\mathbf{c}$) for FPSs. (I did this above for the sake of
convenience, but this is rarely done in the literature.)

\subsubsection{Inverses in commutative rings}

We recall the notion of an \emph{inverse} in a commutative ring:

\begin{definition}
\label{def.commring.inverse}Let $L$ be a commutative ring. Let $a\in L$. Then:
\medskip

\textbf{(a)} An \emph{inverse} (or \emph{multiplicative inverse}) of $a$ means
an element $b\in L$ such that $ab=ba=1$ (where the \textquotedblleft%
$1$\textquotedblright\ means the unity of $L$). \medskip

\textbf{(b)} We say that $a$ is \emph{invertible} in $L$ (or a \emph{unit} of
$L$) if $a$ has an inverse.
\end{definition}

Note that the condition \textquotedblleft$ab=ba=1$\textquotedblright\ in
Definition \ref{def.commring.inverse} \textbf{(a)} can be restated as
\textquotedblleft$ab=1$\textquotedblright, because we automatically have
$ab=ba$ (since $L$ is a commutative ring). I have chosen to write
\textquotedblleft$ab=ba=1$\textquotedblright\ in order to state the definition
in a form that applies verbatim to noncommutative rings as well.

\begin{example}
\textbf{(a)} In the ring $\mathbb{Z}$, the only two invertible elements are
$1$ and $-1$. Each of these two elements is its own inverse. \medskip

\textbf{(b)} In the ring $\mathbb{Q}$, every nonzero element is invertible.
The same holds for the rings $\mathbb{R}$ and $\mathbb{C}$ (and, more
generally, for any field).
\end{example}

Our next goal is to study inverses of FPSs in $K\left[  \left[  x\right]
\right]  $, answering in particular the natural question \textquotedblleft
which elements of $K\left[  \left[  x\right]  \right]  $ have
inverses\textquotedblright. But let us first prove their uniqueness in the
generality of an arbitrary commutative ring:

\begin{theorem}
\label{thm.commring.inverse-uni}Let $L$ be a commutative ring. Let $a\in L$.
Then, there is \textbf{at most one} inverse of $a$.
\end{theorem}

\begin{proof}
Let $b$ and $c$ be two inverses of $a$. We must prove that $b=c$.

Since $b$ is an inverse of $a$, we have $ab=ba=1$. Since $c$ is an inverse of
$a$, we have $ac=ca=1$. Now, we have $b\underbrace{\left(  ac\right)  }%
_{=1}=b\cdot1=b$ and $\underbrace{\left(  ba\right)  }_{=1}c=1\cdot c=c$.
However, because of the \textquotedblleft associativity of
multiplication\textquotedblright\ axiom in Definition \ref{def.alg.commring},
we have $b\left(  ac\right)  =\left(  ba\right)  c$. Hence, $b=b\left(
ac\right)  =\left(  ba\right)  c=c$. This proves Theorem
\ref{thm.commring.inverse-uni}.
\end{proof}

Theorem \ref{thm.commring.inverse-uni} allows us to make the following definition:

\begin{definition}
\label{def.commring.fracs}Let $L$ be a commutative ring. Let $a\in L$. Assume
that $a$ is invertible. Then: \medskip

\textbf{(a)} The inverse of $a$ is called $a^{-1}$. (This notation is
unambiguous, since Theorem \ref{thm.commring.inverse-uni} shows that the
inverse of $a$ is unique.) \medskip

\textbf{(b)} For any $b\in L$, the product $b\cdot a^{-1}$ is called
$\dfrac{b}{a}$ (or $b/a$). \medskip

\textbf{(c)} For any negative integer $n$, we define $a^{n}$ to be $\left(
a^{-1}\right)  ^{-n}$. Thus, the $n$-th power $a^{n}$ is defined for each
$n\in\mathbb{Z}$.
\end{definition}

The following facts are easy to check:

\begin{proposition}
\label{prop.commring.fracs.1}Let $L$ be a commutative ring. Then: \medskip

\textbf{(a)} Any invertible element $a\in L$ satisfies $a^{-1}=1/a$. \medskip

\textbf{(b)} For any invertible elements $a,b\in L$, the element $ab$ is
invertible as well, and satisfies $\left(  ab\right)  ^{-1}=b^{-1}%
a^{-1}=a^{-1}b^{-1}$. \medskip

\textbf{(c)} If $a\in L$ is invertible, and if $n\in\mathbb{Z}$ is arbitrary,
then $a^{-n}=\left(  a^{-1}\right)  ^{n}=\left(  a^{n}\right)  ^{-1}$.
\medskip

\textbf{(d)} Laws of exponents hold for negative exponents as well: If $a$ and
$b$ are invertible elements of $L$, then%
\begin{align*}
a^{n+m}  &  =a^{n}a^{m}\ \ \ \ \ \ \ \ \ \ \text{for all }n,m\in\mathbb{Z};\\
\left(  ab\right)  ^{n}  &  =a^{n}b^{n}\ \ \ \ \ \ \ \ \ \ \text{for all }%
n\in\mathbb{Z};\\
\left(  a^{n}\right)  ^{m}  &  =a^{nm}\ \ \ \ \ \ \ \ \ \ \text{for all
}n,m\in\mathbb{Z}.
\end{align*}


\textbf{(e)} Laws of fractions hold: If $a$ and $c$ are two invertible
elements of $L$, and if $b$ and $d$ are any two elements of $L$, then
\[
\dfrac{b}{a}+\dfrac{d}{c}=\dfrac{bc+ad}{ac}\ \ \ \ \ \ \ \ \ \ \text{and}%
\ \ \ \ \ \ \ \ \ \ \dfrac{b}{a}\cdot\dfrac{d}{c}=\dfrac{bd}{ac}.
\]


\textbf{(f)} Division undoes multiplication: If $a,b,c$ are three elements of
$L$ with $a$ being invertible, then the equality $\dfrac{c}{a}=b$ is
equivalent to $c=ab$.
\end{proposition}

\begin{proof}
Exercise. (See, e.g., \cite[solution to Exercise 4.1.1]{19s} for a proof of
parts \textbf{(c)} and \textbf{(d)} in the special case where $L=\mathbb{C}$;
essentially the same argument works in the general case. The remaining parts
of Proposition \ref{prop.commring.fracs.1} are even easier to check. Note that
parts \textbf{(a)} and \textbf{(c)} as well as the $\left(  ab\right)
^{-1}=b^{-1}a^{-1}$ part of part \textbf{(b)} would hold even if $L$ was a
noncommutative ring.)
\end{proof}

\subsubsection{Inverses in $K\left[  \left[  x\right]  \right]  $}

Now, which FPSs are invertible in the ring $K\left[  \left[  x\right]
\right]  $ ? For example, we know from (\ref{eq.sec.gf.exas.1.1/1-x}) that the
FPS $1-x$ is invertible, with inverse $1+x+x^{2}+x^{3}+\cdots$. On the other
hand, the FPS $x$ is not invertible, since Lemma \ref{lem.fps.xa} shows that
any product of $x$ with an FPS must begin with a $0$ (but the unity of
$K\left[  \left[  x\right]  \right]  $ does not begin with a $0$). (Strictly
speaking, this is only true if the ring $K$ is nontrivial -- i.e., if not all
elements of $K$ are equal. If $K$ is trivial, then $K\left[  \left[  x\right]
\right]  $ is trivial, and thus any FPS in $K\left[  \left[  x\right]
\right]  $ is invertible, but this does not make an interesting statement.)

It turns out that we can characterize invertible FPSs in $K\left[  \left[
x\right]  \right]  $ in a rather simple way:

\begin{proposition}
\label{prop.fps.invertible}Let $a\in K\left[  \left[  x\right]  \right]  $.
Then, the FPS $a$ is invertible in $K\left[  \left[  x\right]  \right]  $ if
and only if its constant term $\left[  x^{0}\right]  a$ is invertible in $K$.
\end{proposition}

\begin{proof}
The statement we are proving is an \textquotedblleft if and only
if\textquotedblright\ statement. We shall prove its \textquotedblleft only
if\textquotedblright\ (i.e., \textquotedblleft$\Longrightarrow$%
\textquotedblright) and its \textquotedblleft if\textquotedblright\ (i.e.,
\textquotedblleft$\Longleftarrow$\textquotedblright) directions separately:

\begin{enumerate}
\item[$\Longrightarrow:$] Assume that $a$ is invertible in $K\left[  \left[
x\right]  \right]  $. That is, $a$ has an inverse $b\in K\left[  \left[
x\right]  \right]  $. Consider this $b$. Since $b$ is an inverse of $a$, we
have $ab=ba=1$ (where \textquotedblleft$1$\textquotedblright\ means
$\underline{1}$ by Convention \ref{conv.fps.const}). However,
(\ref{pf.thm.fps.ring.x0(ab)=}) (applied to $\mathbf{a}=a$ and $\mathbf{b}=b$)
yields $\left[  x^{0}\right]  \left(  ab\right)  =\left[  x^{0}\right]
a\cdot\left[  x^{0}\right]  b$. Comparing this with $\left[  x^{0}\right]
\underbrace{\left(  ab\right)  }_{=1}=\left[  x^{0}\right]  1=1$, we find
$\left[  x^{0}\right]  a\cdot\left[  x^{0}\right]  b=1$. Thus, $\left[
x^{0}\right]  b$ is an inverse of $\left[  x^{0}\right]  a$ in $K$ (since
$\left[  x^{0}\right]  b\cdot\left[  x^{0}\right]  a=\left[  x^{0}\right]
a\cdot\left[  x^{0}\right]  b=1$, so that $\left[  x^{0}\right]  a\cdot\left[
x^{0}\right]  b=\left[  x^{0}\right]  b\cdot\left[  x^{0}\right]  a=1$).
Therefore, $\left[  x^{0}\right]  a$ is invertible in $K$ (with inverse
$\left[  x^{0}\right]  b$). This proves the \textquotedblleft$\Longrightarrow
$\textquotedblright\ direction of Proposition \ref{prop.fps.invertible}.

\item[$\Longleftarrow:$] Assume that $\left[  x^{0}\right]  a$ is invertible
in $K$. Write the FPS $a$ in the form $a=\left(  a_{0},a_{1},a_{2}%
,\ldots\right)  $. Thus, $\left[  x^{0}\right]  a=a_{0}$, so that $a_{0}$ is
invertible in $K$ (since $\left[  x^{0}\right]  a$ is invertible in $K$).
Thus, its inverse $a_{0}^{-1}$ is well-defined.

Now, we want to prove that $a$ is invertible in $K\left[  \left[  x\right]
\right]  $. We thus try to find an inverse of $a$.

We work backwards at first: We assume that $b=\left(  b_{0},b_{1},b_{2}%
,\ldots\right)  \in K\left[  \left[  x\right]  \right]  $ is an inverse for
$a$, and we try to figure out how this inverse looks like.

Since $b$ is an inverse of $a$, we have $ab=\underline{1}=\left(
1,0,0,0,\ldots\right)  $. However, from $a=\left(  a_{0},a_{1},a_{2}%
,\ldots\right)  $ and $b=\left(  b_{0},b_{1},b_{2},\ldots\right)  $, we have%
\begin{align*}
ab  &  =\left(  a_{0},a_{1},a_{2},\ldots\right)  \left(  b_{0},b_{1}%
,b_{2},\ldots\right) \\
&  =\left(  a_{0}b_{0},\ \ a_{0}b_{1}+a_{1}b_{0},\ \ a_{0}b_{2}+a_{1}%
b_{1}+a_{2}b_{0},\ \ a_{0}b_{3}+a_{1}b_{2}+a_{2}b_{1}+a_{3}b_{0}%
,\ \ \ldots\right)
\end{align*}
(by the definition of the product of FPSs). Comparing this with $ab=\left(
1,0,0,0,\ldots\right)  $, we obtain%
\begin{align*}
&  \left(  1,0,0,0,\ldots\right) \\
&  =\left(  a_{0}b_{0},\ \ a_{0}b_{1}+a_{1}b_{0},\ \ a_{0}b_{2}+a_{1}%
b_{1}+a_{2}b_{0},\ \ a_{0}b_{3}+a_{1}b_{2}+a_{2}b_{1}+a_{3}b_{0}%
,\ \ \ldots\right)  .
\end{align*}
This can be rewritten as the following system of equations:%
\begin{equation}
\left\{
\begin{array}
[c]{l}%
1=a_{0}b_{0},\\
0=a_{0}b_{1}+a_{1}b_{0},\\
0=a_{0}b_{2}+a_{1}b_{1}+a_{2}b_{0},\\
0=a_{0}b_{3}+a_{1}b_{2}+a_{2}b_{1}+a_{3}b_{0},\\
\ldots.
\end{array}
\right.  \label{pf.prop.fps.invertible.sys}%
\end{equation}
I claim that this system of equations uniquely determines $\left(  b_{0}%
,b_{1},b_{2},\ldots\right)  $. Indeed, we can solve the first equation
($1=a_{0}b_{0}$) for $b_{0}$, thus obtaining $b_{0}=a_{0}^{-1}$ (since $a_{0}$
is invertible). Having thus found $b_{0}$, we can solve the second equation
($0=a_{0}b_{1}+a_{1}b_{0}$) for $b_{1}$, thus obtaining $b_{1}=-a_{0}%
^{-1}\left(  a_{1}b_{0}\right)  $ (again because $a_{0}$ is invertible).
Having thus found both $b_{0}$ and $b_{1}$, we can solve the third equation
($0=a_{0}b_{2}+a_{1}b_{1}+a_{2}b_{0}$) for $b_{2}$, thus obtaining
$b_{2}=-a_{0}^{-1}\left(  a_{1}b_{1}+a_{2}b_{0}\right)  $. Proceeding like
this, we obtain recursive expressions for all coefficients $b_{0},b_{1}%
,b_{2},\ldots$ of $b$, namely%
\begin{equation}
\left\{
\begin{array}
[c]{l}%
b_{0}=a_{0}^{-1},\\
b_{1}=-a_{0}^{-1}\left(  a_{1}b_{0}\right)  ,\\
b_{2}=-a_{0}^{-1}\left(  a_{1}b_{1}+a_{2}b_{0}\right)  ,\\
b_{3}=-a_{0}^{-1}\left(  a_{1}b_{2}+a_{2}b_{1}+a_{3}b_{0}\right)  ,\\
\ldots.
\end{array}
\right.  \label{pf.prop.fps.invertible.sol}%
\end{equation}
(This procedure for solving systems of linear equations is well-known from
linear algebra -- it is a form of Gaussian elimination, but a particularly
simple one because our system is triangular with invertible coefficients on
the diagonal. The only complication is that it has infinitely many variables
and infinitely many equations.)

So we have shown that if $b$ is an inverse of $a$, then the entries $b_{i}$ of
the FPS $b$ are given recursively by (\ref{pf.prop.fps.invertible.sol}). This
yields that $b$ is unique; alas, this is not what we want to prove. Instead,
we want to prove that $b$ exists.

Fortunately, we can achieve this by simply turning our above argument around:
Forget that we fixed $b$. Instead, we define a sequence $\left(  b_{0}%
,b_{1},b_{2},\ldots\right)  $ of elements of $K$ recursively by
(\ref{pf.prop.fps.invertible.sol}), and we define the FPS $b=\left(
b_{0},b_{1},b_{2},\ldots\right)  \in K\left[  \left[  x\right]  \right]
$.\ Then, the equalities (\ref{pf.prop.fps.invertible.sys}) hold (because they
are just equivalent restatements of the equalities
(\ref{pf.prop.fps.invertible.sol})). In other words, we have%
\begin{align*}
&  \left(  1,0,0,0,\ldots\right) \\
&  =\left(  a_{0}b_{0},\ \ a_{0}b_{1}+a_{1}b_{0},\ \ a_{0}b_{2}+a_{1}%
b_{1}+a_{2}b_{0},\ \ a_{0}b_{3}+a_{1}b_{2}+a_{2}b_{1}+a_{3}b_{0}%
,\ \ \ldots\right)  .
\end{align*}
However, as before, we can show that%
\[
ab=\left(  a_{0}b_{0},\ \ a_{0}b_{1}+a_{1}b_{0},\ \ a_{0}b_{2}+a_{1}%
b_{1}+a_{2}b_{0},\ \ a_{0}b_{3}+a_{1}b_{2}+a_{2}b_{1}+a_{3}b_{0}%
,\ \ \ldots\right)  .
\]
Comparing these two equalities, we find $ab=\left(  1,0,0,0,\ldots\right)
=\underline{1}$. Thus, $ba=ab=\underline{1}$, so that $ab=ba=\underline{1}$.
This shows that $b$ is an inverse of $a$, so that $a$ is invertible. This
proves the \textquotedblleft$\Longleftarrow$\textquotedblright\ direction of
Proposition \ref{prop.fps.invertible}.
\end{enumerate}
\end{proof}

We note a particularly simple corollary of Proposition
\ref{prop.fps.invertible} when $K$ is a field:

\begin{corollary}
\label{cor.fps.invertible.field}Assume that $K$ is a field. Let $a\in K\left[
\left[  x\right]  \right]  $. Then, the FPS $a$ is invertible in $K\left[
\left[  x\right]  \right]  $ if and only if $\left[  x^{0}\right]  a\neq0$.
\end{corollary}

\begin{proof}
An element of $K$ is invertible in $K$ if and only if it is nonzero (since $K$
is a field). Hence, Corollary \ref{cor.fps.invertible.field} follows from
Proposition \ref{prop.fps.invertible}.
\end{proof}

\subsubsection{Newton's binomial formula}

Let us now return to considering specific FPSs. We have already seen that the
FPS $1-x$ is invertible, with inverse $1+x+x^{2}+x^{3}+\cdots$. We shall now
show an analogous result for the FPS $1+x$. Its invertibility follows from
Proposition \ref{prop.fps.invertible}, but it is better to derive it by hand,
as this also gives a formula for the inverse:

\begin{proposition}
\label{prop.fps.invertible.1+x}The FPS $1+x\in K\left[  \left[  x\right]
\right]  $ is invertible, and its inverse is%
\[
\left(  1+x\right)  ^{-1}=1-x+x^{2}-x^{3}+x^{4}-x^{5}\pm\cdots=\sum
_{n\in\mathbb{N}}\left(  -1\right)  ^{n}x^{n}.
\]

\end{proposition}

\begin{proof}
[First proof of Proposition \ref{prop.fps.invertible.1+x}.]We have%
\begin{align*}
&  \left(  1+x\right)  \cdot\left(  1-x+x^{2}-x^{3}+x^{4}-x^{5}\pm
\cdots\right) \\
&  =\underbrace{1\cdot\left(  1-x+x^{2}-x^{3}+x^{4}-x^{5}\pm\cdots\right)
}_{=1-x+x^{2}-x^{3}+x^{4}-x^{5}\pm\cdots}+\underbrace{x\cdot\left(
1-x+x^{2}-x^{3}+x^{4}-x^{5}\pm\cdots\right)  }_{=x-x^{2}+x^{3}-x^{4}%
+x^{5}-x^{6}\pm\cdots}\\
&  =\left(  1-x+x^{2}-x^{3}+x^{4}-x^{5}\pm\cdots\right)  +\left(
x-x^{2}+x^{3}-x^{4}+x^{5}-x^{6}\pm\cdots\right) \\
&  =1
\end{align*}
(since all powers of $x$ other than $1$ cancel out). This shows that
$1-x+x^{2}-x^{3}+x^{4}-x^{5}\pm\cdots$ is an inverse of $1+x$ (since $K\left[
\left[  x\right]  \right]  $ is a commutative ring). Thus, $1+x$ is
invertible, and its inverse is $\left(  1+x\right)  ^{-1}=1-x+x^{2}%
-x^{3}+x^{4}-x^{5}\pm\cdots=\sum_{n\in\mathbb{N}}\left(  -1\right)  ^{n}x^{n}%
$. This proves Proposition \ref{prop.fps.invertible.1+x}.
\end{proof}

\begin{proof}
[Second proof of Proposition \ref{prop.fps.invertible.1+x}.]We have%
\begin{align*}
&  \left(  1+x\right)  \cdot\left(  1-x+x^{2}-x^{3}+x^{4}-x^{5}\pm
\cdots\right) \\
&  =\underbrace{\left(  1+x\right)  \cdot1}_{=1+x}-\underbrace{\left(
1+x\right)  \cdot x}_{=x+x^{2}}+\underbrace{\left(  1+x\right)  \cdot x^{2}%
}_{=x^{2}+x^{3}}-\underbrace{\left(  1+x\right)  \cdot x^{3}}_{=x^{3}+x^{4}%
}+\underbrace{\left(  1+x\right)  \cdot x^{4}}_{=x^{4}+x^{5}}%
-\underbrace{\left(  1+x\right)  \cdot x^{5}}_{=x^{5}+x^{6}}\pm\cdots\\
&  =\left(  1+x\right)  -\left(  x+x^{2}\right)  +\left(  x^{2}+x^{3}\right)
-\left(  x^{3}+x^{4}\right)  +\left(  x^{4}+x^{5}\right)  -\left(  x^{5}%
+x^{6}\right)  \pm\cdots\\
&  =1
\end{align*}
(since we have a telescoping sum in front of us, in which all powers of $x$
other than $1$ cancel out). This shows that $1-x+x^{2}-x^{3}+x^{4}-x^{5}%
\pm\cdots$ is an inverse of $1+x$ (since $K\left[  \left[  x\right]  \right]
$ is a commutative ring). Thus, $1+x$ is invertible, and its inverse is
$\left(  1+x\right)  ^{-1}=1-x+x^{2}-x^{3}+x^{4}-x^{5}\pm\cdots=\sum
_{n\in\mathbb{N}}\left(  -1\right)  ^{n}x^{n}$. This proves Proposition
\ref{prop.fps.invertible.1+x}.
\end{proof}

Proposition \ref{prop.fps.invertible.1+x} shows that the FPS $1+x$ is
invertible; thus, its powers $\left(  1+x\right)  ^{n}$ are defined for all
$n\in\mathbb{Z}$ (by Definition \ref{def.commring.fracs} \textbf{(c)}). The
following formula -- known as \emph{Newton's binomial theorem}\footnote{or
\emph{Newton's binomial formula}} -- describes these powers explicitly:

\begin{theorem}
\label{thm.fps.newton-binom}For each $n\in\mathbb{Z}$, we have%
\[
\left(  1+x\right)  ^{n}=\sum_{k\in\mathbb{N}}\dbinom{n}{k}x^{k}.
\]

\end{theorem}

Note that the sum $\sum_{k\in\mathbb{N}}\dbinom{n}{k}x^{k}$ is summable for
each $n\in\mathbb{N}$ (indeed, it equals the FPS $\left(  \dbinom{n}%
{0},\dbinom{n}{1},\dbinom{n}{2},\ldots\right)  $). If $n\in\mathbb{N}$, then
it is essentially finite.

The reader may want to check that the particular case $n=-1$ of Theorem
\ref{thm.fps.newton-binom} agrees with Proposition
\ref{prop.fps.invertible.1+x}. (Recall Example \ref{exa.binom.-1choosek}!)

Of course, Theorem \ref{thm.fps.newton-binom} should look familiar -- an
identical-looking formula appears in real analysis under the same name.
However, the result in real analysis is concerned with infinite sums of real
numbers, while our Theorem \ref{thm.fps.newton-binom} is an identity between
FPSs over an arbitrary commutative ring. Thus, the two facts are not the same.

We will prove Theorem \ref{thm.fps.newton-binom} in a somewhat roundabout way,
since this gives us an opportunity to establish some auxiliary results that
are of separate interest (and usefulness). The first of these auxiliary
results is a fundamental property of binomial coefficients, known as the
\emph{upper negation formula} (see, e.g., \cite[Proposition 1.3.7]{19fco}):

\begin{theorem}
\label{thm.binom.upneg-n}Let $n\in\mathbb{C}$ and $k\in\mathbb{Z}$. Then,%
\[
\dbinom{-n}{k}=\left(  -1\right)  ^{k}\dbinom{k+n-1}{k}.
\]

\end{theorem}

\begin{proof}
[Proof of Theorem \ref{thm.binom.upneg-n} (sketched).]If $k<0$, then this is
trivial because both $\dbinom{-n}{k}$ and $\dbinom{k+n-1}{k}$ are $0$ (by
(\ref{eq.def.binom.binom.eq})). Thus, we WLOG assume that $k\geq0$. Hence,
$k\in\mathbb{N}$. Thus, (\ref{eq.def.binom.binom.eq}) yields%
\begin{align*}
\dbinom{-n}{k}  &  =\dfrac{\left(  -n\right)  \left(  -n-1\right)  \left(
-n-2\right)  \cdots\left(  -n-k+1\right)  }{k!}\ \ \ \ \ \ \ \ \ \ \text{and}%
\\
\dbinom{k+n-1}{k}  &  =\dfrac{\left(  k+n-1\right)  \left(  k+n-2\right)
\left(  k+n-3\right)  \cdots\left(  k+n-k\right)  }{k!}.
\end{align*}
A moment of thought reveals that the right hand sides of these two equalities
are equal up to a factor of $\left(  -1\right)  ^{k}$. Thus, so are the left
hand sides. In other words, $\dbinom{-n}{k}=\left(  -1\right)  ^{k}%
\dbinom{k+n-1}{k}$. This proves Theorem \ref{thm.binom.upneg-n}.
\end{proof}

(Quick exercise: Rederive Example \ref{exa.binom.-1choosek} from Theorem
\ref{thm.binom.upneg-n}.)

Next, we show a formula for negative powers of $1+x$:

\begin{proposition}
\label{prop.fps.anti-newton-binom}For each $n\in\mathbb{N}$, we have%
\[
\left(  1+x\right)  ^{-n}=\sum_{k\in\mathbb{N}}\left(  -1\right)  ^{k}%
\dbinom{n+k-1}{k}x^{k}.
\]

\end{proposition}

\begin{proof}
[Proof of Proposition \ref{prop.fps.anti-newton-binom}.]We proceed by
induction on $n$:

\textit{Induction base:} Comparing
\[
\left(  1+x\right)  ^{-0}=\left(  1+x\right)  ^{0}=\underline{1}=\left(
1,0,0,0,\ldots\right)  =1
\]
with%
\begin{align*}
&  \sum_{k\in\mathbb{N}}\left(  -1\right)  ^{k}\dbinom{0+k-1}{k}x^{k}\\
&  =\underbrace{\left(  -1\right)  ^{0}}_{=1}\underbrace{\dbinom{0+0-1}{0}%
}_{=1}\underbrace{x^{0}}_{=1}+\sum_{\substack{k\in\mathbb{N};\\k>0}}\left(
-1\right)  ^{k}\underbrace{\dbinom{0+k-1}{k}}_{\substack{=\dbinom{k-1}%
{k}=0\\\text{(by Proposition \ref{prop.binom.0}}\\\text{(applied to
}m=k-1\text{ and }n=k\text{),}\\\text{since }k-1<k\text{ and }k-1\in
\mathbb{N}\text{)}}}x^{k}\\
&  \ \ \ \ \ \ \ \ \ \ \ \ \ \ \ \ \ \ \ \ \left(  \text{here, we have split
off the addend for }k=0\text{ from the sum}\right) \\
&  =1+\underbrace{\sum_{\substack{k\in\mathbb{N};\\k>0}}\left(  -1\right)
^{k}0x^{k}}_{=0}=1,
\end{align*}
we obtain $\left(  1+x\right)  ^{-0}=\sum_{k\in\mathbb{N}}\left(  -1\right)
^{k}\dbinom{0+k-1}{k}x^{k}$. In other words, Proposition
\ref{prop.fps.anti-newton-binom} holds for $n=0$.

\textit{Induction step:} Let $j\in\mathbb{N}$. Assume that Proposition
\ref{prop.fps.anti-newton-binom} holds for $n=j$. We must prove that
Proposition \ref{prop.fps.anti-newton-binom} holds for $n=j+1$.

We have assumed that Proposition \ref{prop.fps.anti-newton-binom} holds for
$n=j$. In other words, we have%
\begin{equation}
\left(  1+x\right)  ^{-j}=\sum_{k\in\mathbb{N}}\left(  -1\right)  ^{k}%
\dbinom{j+k-1}{k}x^{k}. \label{pf.prop.fps.anti-newton-binom.IH}%
\end{equation}


Now, we want to prove that Proposition \ref{prop.fps.anti-newton-binom} holds
for $n=j+1$. In other words, we want to prove that%
\[
\left(  1+x\right)  ^{-\left(  j+1\right)  }=\sum_{k\in\mathbb{N}}\left(
-1\right)  ^{k}\dbinom{\left(  j+1\right)  +k-1}{k}x^{k}.
\]
In view of $\left(  1+x\right)  ^{-\left(  j+1\right)  }=\left(  1+x\right)
^{-j}\cdot\left(  1+x\right)  ^{-1}$ and $\left(  j+1\right)  +k-1=j+k$, this
equality can be rewritten as%
\[
\left(  1+x\right)  ^{-j}\cdot\left(  1+x\right)  ^{-1}=\sum_{k\in\mathbb{N}%
}\left(  -1\right)  ^{k}\dbinom{j+k}{k}x^{k}.
\]
Since $1+x$ is invertible, we can equivalently transform this equality by
multiplying both sides with $1+x$; thus, it becomes%
\[
\left(  1+x\right)  ^{-j}=\left(  \sum_{k\in\mathbb{N}}\left(  -1\right)
^{k}\dbinom{j+k}{k}x^{k}\right)  \cdot\left(  1+x\right)  .
\]
So this is the equality we must prove.

We do this by simplifying its right hand side:%
\begin{align*}
&  \left(  \sum_{k\in\mathbb{N}}\left(  -1\right)  ^{k}\dbinom{j+k}{k}%
x^{k}\right)  \cdot\left(  1+x\right) \\
&  =\sum_{k\in\mathbb{N}}\underbrace{\left(  -1\right)  ^{k}\dbinom{j+k}%
{k}x^{k}\left(  1+x\right)  }_{\substack{=\left(  -1\right)  ^{k}\dbinom
{j+k}{k}\left(  x^{k}+x^{k+1}\right)  \\=\left(  -1\right)  ^{k}\dbinom
{j+k}{k}x^{k}+\left(  -1\right)  ^{k}\dbinom{j+k}{k}x^{k+1}}}\\
&  =\sum_{k\in\mathbb{N}}\left(  \left(  -1\right)  ^{k}\dbinom{j+k}{k}%
x^{k}+\left(  -1\right)  ^{k}\dbinom{j+k}{k}x^{k+1}\right) \\
&  =\sum_{k\in\mathbb{N}}\left(  -1\right)  ^{k}\underbrace{\dbinom{j+k}{k}%
}_{\substack{=\dbinom{j+k-1}{k-1}+\dbinom{j+k-1}{k}\\\text{(by Proposition
\ref{prop.binom.rec},}\\\text{applied to }m=j+k\text{ and }n=k\text{)}}%
}x^{k}+\underbrace{\sum_{k\in\mathbb{N}}\left(  -1\right)  ^{k}\dbinom{j+k}%
{k}x^{k+1}}_{\substack{=\sum_{k\geq1}\left(  -1\right)  ^{k-1}\dbinom
{j+\left(  k-1\right)  }{k-1}x^{k}\\\text{(here, we have substituted
}k-1\text{ for }k\\\text{in the sum)}}}\\
&  =\underbrace{\sum_{k\in\mathbb{N}}\left(  -1\right)  ^{k}\left(
\dbinom{j+k-1}{k-1}+\dbinom{j+k-1}{k}\right)  x^{k}}_{=\sum_{k\in\mathbb{N}%
}\left(  -1\right)  ^{k}\dbinom{j+k-1}{k-1}x^{k}+\sum_{k\in\mathbb{N}}\left(
-1\right)  ^{k}\dbinom{j+k-1}{k}x^{k}}+\sum_{k\geq1}\underbrace{\left(
-1\right)  ^{k-1}}_{=-\left(  -1\right)  ^{k}}\underbrace{\dbinom{j+\left(
k-1\right)  }{k-1}}_{=\dbinom{j+k-1}{k-1}}x^{k}\\
&  =\sum_{k\in\mathbb{N}}\left(  -1\right)  ^{k}\dbinom{j+k-1}{k-1}x^{k}%
+\sum_{k\in\mathbb{N}}\left(  -1\right)  ^{k}\dbinom{j+k-1}{k}x^{k}%
+\sum_{k\geq1}\left(  -\left(  -1\right)  ^{k}\right)  \dbinom{j+k-1}%
{k-1}x^{k}\\
&  =\sum_{k\in\mathbb{N}}\left(  -1\right)  ^{k}\dbinom{j+k-1}{k-1}x^{k}%
+\sum_{k\in\mathbb{N}}\left(  -1\right)  ^{k}\dbinom{j+k-1}{k}x^{k}%
-\sum_{k\geq1}\left(  -1\right)  ^{k}\dbinom{j+k-1}{k-1}x^{k}\\
&  =\underbrace{\left(  \sum_{k\in\mathbb{N}}\left(  -1\right)  ^{k}%
\dbinom{j+k-1}{k-1}x^{k}-\sum_{k\geq1}\left(  -1\right)  ^{k}\dbinom
{j+k-1}{k-1}x^{k}\right)  }_{\substack{=\left(  -1\right)  ^{0}\dbinom
{j+0-1}{0-1}x^{0}\\\text{(since the two sums differ only in their }k=0\text{
addend)}}}+\underbrace{\sum_{k\in\mathbb{N}}\left(  -1\right)  ^{k}%
\dbinom{j+k-1}{k}x^{k}}_{\substack{=\left(  1+x\right)  ^{-j}\\\text{(by
(\ref{pf.prop.fps.anti-newton-binom.IH}))}}}\\
&  =\left(  -1\right)  ^{0}\underbrace{\dbinom{j+0-1}{0-1}}%
_{\substack{=0\\\text{(by (\ref{eq.def.binom.binom.eq}), since }%
0-1\notin\mathbb{N}\text{)}}}x^{0}+\left(  1+x\right)  ^{-j}=\left(
1+x\right)  ^{-j}.
\end{align*}
Multiplying both sides of this equality by $\left(  1+x\right)  ^{-1}$, we
obtain%
\[
\sum_{k\in\mathbb{N}}\left(  -1\right)  ^{k}\dbinom{j+k}{k}x^{k}=\left(
1+x\right)  ^{-j}\cdot\left(  1+x\right)  ^{-1}=\left(  1+x\right)  ^{\left(
-j\right)  +\left(  -1\right)  }=\left(  1+x\right)  ^{-\left(  j+1\right)
},
\]
so that
\[
\left(  1+x\right)  ^{-\left(  j+1\right)  }=\sum_{k\in\mathbb{N}}\left(
-1\right)  ^{k}\underbrace{\dbinom{j+k}{k}}_{=\dbinom{\left(  j+1\right)
+k-1}{k}}x^{k}=\sum_{k\in\mathbb{N}}\left(  -1\right)  ^{k}\dbinom{\left(
j+1\right)  +k-1}{k}x^{k}.
\]
In other words, Proposition \ref{prop.fps.anti-newton-binom} holds for
$n=j+1$. This completes the induction step. Thus, Proposition
\ref{prop.fps.anti-newton-binom} is proved.
\end{proof}

We can rewrite Proposition \ref{prop.fps.anti-newton-binom} using negative
binomial coefficients:

\begin{corollary}
\label{cor.fps.anti-newton-binom-2}For each $n\in\mathbb{N}$, we have%
\[
\left(  1+x\right)  ^{-n}=\sum_{k\in\mathbb{N}}\dbinom{-n}{k}x^{k}.
\]

\end{corollary}

\begin{proof}
[Proof of Corollary \ref{cor.fps.anti-newton-binom-2}.]Proposition
\ref{prop.fps.anti-newton-binom} yields%
\[
\left(  1+x\right)  ^{-n}=\sum_{k\in\mathbb{N}}\underbrace{\left(  -1\right)
^{k}\dbinom{n+k-1}{k}}_{\substack{=\left(  -1\right)  ^{k}\dbinom{k+n-1}%
{k}\\=\dbinom{-n}{k}\\\text{(by Theorem \ref{thm.binom.upneg-n})}}}x^{k}%
=\sum_{k\in\mathbb{N}}\dbinom{-n}{k}x^{k}.
\]
This proves Corollary \ref{cor.fps.anti-newton-binom-2}.
\end{proof}

We can now easily prove Newton's binomial formula:

\begin{proof}
[Proof of Theorem \ref{thm.fps.newton-binom}.]Let $n\in\mathbb{Z}$. We must
prove that $\left(  1+x\right)  ^{n}=\sum_{k\in\mathbb{N}}\dbinom{n}{k}x^{k}$.

If $n\in\mathbb{N}$, then this follows by comparing%
\begin{align*}
\left(  1+x\right)  ^{n}  &  =\left(  x+1\right)  ^{n}=\sum_{k=0}^{n}%
\dbinom{n}{k}x^{k}\underbrace{1^{n-k}}_{=1}\ \ \ \ \ \ \ \ \ \ \left(
\text{by the binomial theorem}\right) \\
&  =\sum_{k=0}^{n}\dbinom{n}{k}x^{k}%
\end{align*}
with%
\begin{align*}
\sum_{k\in\mathbb{N}}\dbinom{n}{k}x^{k}  &  =\sum_{k=0}^{n}\dbinom{n}{k}%
x^{k}+\sum_{k>n}\underbrace{\dbinom{n}{k}}_{\substack{=0\\\text{(by
Proposition \ref{prop.binom.0},}\\\text{since }n<k\text{)}}}x^{k}=\sum
_{k=0}^{n}\dbinom{n}{k}x^{k}+\underbrace{\sum_{k>n}0x^{k}}_{=0}\\
&  =\sum_{k=0}^{n}\dbinom{n}{k}x^{k}.
\end{align*}
Hence, for the rest of this proof, we WLOG assume that $n\notin\mathbb{N}$.

Hence, $n$ is a negative integer, so that $-n\in\mathbb{N}$. Thus, Corollary
\ref{cor.fps.anti-newton-binom-2} (applied to $-n$ instead of $n$) yields%
\[
\left(  1+x\right)  ^{-\left(  -n\right)  }=\sum_{k\in\mathbb{N}}%
\dbinom{-\left(  -n\right)  }{k}x^{k}.
\]
Since $-\left(  -n\right)  =n$, this rewrites as $\left(  1+x\right)
^{n}=\sum_{k\in\mathbb{N}}\dbinom{n}{k}x^{k}$. Thus, Theorem
\ref{thm.fps.newton-binom} is proven.
\end{proof}

We thus have a formula for $\left(  1+x\right)  ^{n}$ for each
\textbf{integer} $n$. We don't yet have such a formula for $\left(
1+x\right)  ^{1/2}$ (nor do we have a proper definition of $\left(
1+x\right)  ^{1/2}$), but this was clearly a step forward.

\subsubsection{Dividing by $x$}

Let us see how this all helps us justify our arguments in Section
\ref{sec.gf.exas}. Proposition \ref{prop.fps.invertible} justifies the
fractions that appear in (\ref{eq.sec.gf.exas.1.Fx=2}), but it does not
justify dividing by the FPS $2x$ in (\ref{eq.sec.gf.exas.2.-}), since the
constant term $\left[  x^{0}\right]  \left(  2x\right)  $ is surely not
invertible. And indeed, the FPS $2x$ is not invertible; the fraction
$\dfrac{1}{2x}$ is not a well-defined FPS.

However, it is easy to see directly which FPSs can be divided by $x$ (and thus
by $2x$, if $K=\mathbb{Q}$), and what it means to divide them by $x$. In fact,
Lemma \ref{lem.fps.xa} shows that multiplying an FPS by $x$ means moving all
its entries by one position to the right, and putting a $0$ into the newly
vacated starting position. Thus, it is rather clear what dividing by $x$
should be:

\begin{definition}
\label{def.fps.div-by-x}Let $a=\left(  a_{0},a_{1},a_{2},\ldots\right)  $ be
an FPS whose constant term $a_{0}$ is $0$. Then, $\dfrac{a}{x}$ is defined to
be the FPS $\left(  a_{1},a_{2},a_{3},\ldots\right)  $.
\end{definition}

The following is almost trivial:

\begin{proposition}
\label{prop.fps.div-by-x-inverts}Let $a\in K\left[  \left[  x\right]  \right]
$ and $b\in K\left[  \left[  x\right]  \right]  $ be two FPSs. Then, $a=xb$ if
and only if $\left(  \left[  x^{0}\right]  a=0\text{ and }b=\dfrac{a}%
{x}\right)  $.
\end{proposition}

\begin{proof}
Exercise.
\end{proof}

Having defined $\dfrac{a}{x}$ in Definition \ref{def.fps.div-by-x} (when $a$
has constant term $0$), we can also define $\dfrac{a}{2x}$ when $2$ is
invertible in $K$ (just set $\dfrac{a}{2x}=\dfrac{1}{2}\cdot\dfrac{a}{x}$).
Thus, the fraction $\dfrac{1\pm\sqrt{1-4x}}{2x}$ in (\ref{eq.sec.gf.exas.2.-})
makes sense when the $\pm$ sign is a $-$ sign (but not when it is a $+$ sign),
at least if we interpret the square root $\sqrt{1-4x}$ as $\sum_{k\geq
0}\dbinom{1/2}{k}\left(  -4x\right)  ^{k}$.

\subsubsection{A few lemmas}

Let us use this occasion to state two simple lemmas (vaguely related to
Definition \ref{def.fps.div-by-x}) that will be used later on:

\begin{lemma}
\label{lem.fps.g=xh}Let $a\in K\left[  \left[  x\right]  \right]  $ be an FPS
with $\left[  x^{0}\right]  a=0$. Then, there exists an $h\in K\left[  \left[
x\right]  \right]  $ such that $a=xh$.
\end{lemma}

\begin{proof}
[Proof of Lemma \ref{lem.fps.g=xh}.]Write the FPS $a$ in the form $a=\left(
a_{0},a_{1},a_{2},\ldots\right)  $. Thus, $a_{0}=\left[  x^{0}\right]  a=0$.
Hence, the FPS $\dfrac{a}{x}$ is well-defined. Moreover, it is easy to see
that $a=x\cdot\dfrac{a}{x}$\ \ \ \ \footnote{\textit{Proof.} We have
$\dfrac{a}{x}=\left(  a_{1},a_{2},a_{3},\ldots\right)  $ (by Definition
\ref{def.fps.div-by-x}). Thus, Lemma \ref{lem.fps.xa} (applied to $\dfrac
{a}{x}$ and $a_{i-1}$ instead of $\mathbf{a}$ and $a_{i}$) yields%
\begin{align*}
x\cdot\dfrac{a}{x}  &  =\left(  0,a_{1},a_{2},a_{3},\ldots\right)  =\left(
a_{0},a_{1},a_{2},a_{3},\ldots\right)  \ \ \ \ \ \ \ \ \ \ \left(  \text{since
}0=a_{0}\right) \\
&  =\left(  a_{0},a_{1},a_{2},\ldots\right)  =a.
\end{align*}
Thus, $a=x\cdot\dfrac{a}{x}$.}. Hence, there exists an $h\in K\left[  \left[
x\right]  \right]  $ such that $a=xh$ (namely, $h=\dfrac{a}{x}$). This proves
Lemma \ref{lem.fps.g=xh}.
\end{proof}

\begin{lemma}
\label{lem.fps.first-n-coeffs-of-xna}Let $k\in\mathbb{N}$. Let $a\in K\left[
\left[  x\right]  \right]  $ be any FPS. Then, the first $k$ coefficients of
the FPS $x^{k}a$ are $0$.
\end{lemma}

\begin{proof}
[Proof of Lemma \ref{lem.fps.first-n-coeffs-of-xna}.]We must show that
$\left[  x^{m}\right]  \left(  x^{k}a\right)  =0$ for any nonnegative integer
$m<k$. But we can do this directly: If $m$ is a nonnegative integer such that
$m<k$, then (\ref{pf.thm.fps.ring.xn(ab)=2}) (applied to $x^{k}$, $a$ and $m$
instead of $\mathbf{a}$, $\mathbf{b}$ and $n$) yields%
\[
\left[  x^{m}\right]  \left(  x^{k}a\right)  =\sum_{i=0}^{m}%
\underbrace{\left[  x^{i}\right]  \left(  x^{k}\right)  }%
_{\substack{=0\\\text{(since }i\leq m<k\\\text{and thus }i\neq k\text{)}%
}}\cdot\left[  x^{m-i}\right]  a=\sum_{i=0}^{m}0\cdot\left[  x^{m-i}\right]
a=0,
\]
which is exactly what we wanted to show. Thus, Lemma
\ref{lem.fps.first-n-coeffs-of-xna} is proved.

(Alternatively, we could prove Lemma \ref{lem.fps.first-n-coeffs-of-xna} by
writing $a$ in the form $a=\left(  a_{0},a_{1},a_{2},\ldots\right)  $ and
observing that $x^{k}a=\left(  \underbrace{0,0,\ldots,0}_{k\text{ times}%
},a_{0},a_{1},a_{2},\ldots\right)  $. This follows by applying Lemma
\ref{lem.fps.xa} a total of $k$ times, or more formally by induction on $k$.)
\end{proof}

Lemma \ref{lem.fps.first-n-coeffs-of-xna} has a converse; here is a statement
that combines it with this converse:

\begin{lemma}
\label{lem.fps.muls-of-xn}Let $k\in\mathbb{N}$. Let $f\in K\left[  \left[
x\right]  \right]  $ be any FPS. Then, the first $k$ coefficients of the FPS
$f$ are $0$ if and only if $f$ is a multiple of $x^{k}$.
\end{lemma}

Here, we use the following notation:

\begin{definition}
Let $g\in K\left[  \left[  x\right]  \right]  $ be an FPS. Then, a
\emph{multiple} of $g$ means an FPS of the form $ga$ with $a\in K\left[
\left[  x\right]  \right]  $.
\end{definition}

(This is just a particular case of the usual concept of multiples in a
commutative ring.)

\begin{proof}
[Proof of Lemma \ref{lem.fps.muls-of-xn}.]The statement we are proving is an
\textquotedblleft if and only if\textquotedblright\ statement. We shall prove
its \textquotedblleft only if\textquotedblright\ (i.e., \textquotedblleft%
$\Longrightarrow$\textquotedblright) and its \textquotedblleft
if\textquotedblright\ (i.e., \textquotedblleft$\Longleftarrow$%
\textquotedblright) directions separately:

\begin{enumerate}
\item[$\Longrightarrow:$] Assume that the first $k$ coefficients of the FPS
$f$ are $0$. We must show that $f$ is a multiple of $x^{k}$.

Write $f$ as $f=\left(  f_{0},f_{1},f_{2},\ldots\right)  $. Then, the first
$k$ coefficients of the FPS $f$ are $f_{0},f_{1},\ldots,f_{k-1}$. Hence, these
$k$ coefficients $f_{0},f_{1},\ldots,f_{k-1}$ are $0$ (since we have assumed
that the first $k$ coefficients of the FPS $f$ are $0$). In other words,
$f_{n}=0$ for each $n\in\left\{  0,1,\ldots,k-1\right\}  $. Hence, $\sum
_{n=0}^{k-1}\underbrace{f_{n}}_{=0}x^{n}=\sum_{n=0}^{k-1}0x^{n}=0$.

Now,
\begin{align*}
f  &  =\left(  f_{0},f_{1},f_{2},\ldots\right)  =\sum_{n\in\mathbb{N}}%
f_{n}x^{n}=\underbrace{\sum_{n=0}^{k-1}f_{n}x^{n}}_{=0}+\sum_{n=k}^{\infty
}f_{n}\underbrace{x^{n}}_{\substack{=x^{k}x^{n-k}\\\text{(since }n\geq
k\text{)}}}=\sum_{n=k}^{\infty}f_{n}x^{k}x^{n-k}\\
&  =x^{k}\sum_{n=k}^{\infty}f_{n}x^{n-k}.
\end{align*}
In other words, $f=x^{k}a$ for $a=\sum_{n=k}^{\infty}f_{n}x^{n-k}$. This shows
that $f$ is a multiple of $x^{k}$. Thus, the \textquotedblleft$\Longrightarrow
$\textquotedblright\ direction of Lemma \ref{lem.fps.muls-of-xn} is proved.

\item[$\Longleftarrow:$] Assume that $f$ is a multiple of $x^{k}$. In other
words, $f=x^{k}a$ for some $a\in K\left[  \left[  x\right]  \right]  $.
Consider this $a$. Now, Lemma \ref{lem.fps.first-n-coeffs-of-xna} yields that
the first $k$ coefficients of the FPS $x^{k}a$ are $0$. In other words, the
first $k$ coefficients of the FPS $f$ are $0$ (since $f=x^{k}a$). This proves
the \textquotedblleft$\Longleftarrow$\textquotedblright\ direction of Lemma
\ref{lem.fps.muls-of-xn}.
\end{enumerate}

The proof of Lemma \ref{lem.fps.muls-of-xn} is now complete, as both
directions have been proved.
\end{proof}

Another lemma that will prove its usefulness much later concerns FPSs that are
equal up until a certain coefficient. It says that if $f$ and $g$ are two FPSs
whose first $n+1$ coefficients agree (for a certain $n\in\mathbb{N}$), then
the same is true of the FPSs $af$ and $ag$ whenever $a$ is any further FPS. In
more details:

\begin{lemma}
\label{lem.fps.prod.irlv.fg}Let $a,f,g\in K\left[  \left[  x\right]  \right]
$ be three FPSs. Let $n\in\mathbb{N}$. Assume that
\begin{equation}
\left[  x^{m}\right]  f=\left[  x^{m}\right]  g\ \ \ \ \ \ \ \ \ \ \text{for
each }m\in\left\{  0,1,\ldots,n\right\}  . \label{eq.lem.fps.prod.irlv.fg.ass}%
\end{equation}
Then,
\[
\left[  x^{m}\right]  \left(  af\right)  =\left[  x^{m}\right]  \left(
ag\right)  \ \ \ \ \ \ \ \ \ \ \text{for each }m\in\left\{  0,1,\ldots
,n\right\}  .
\]

\end{lemma}

\begin{proof}
[Proof of Lemma \ref{lem.fps.prod.irlv.fg}.]Let $m\in\left\{  0,1,\ldots
,n\right\}  $. Then, $m\leq n$. Hence, each $j\in\left\{  0,1,\ldots
,m\right\}  $ satisfies $j\leq m\leq n$ and thus $j\in\left\{  0,1,\ldots
,n\right\}  $ and therefore%
\begin{equation}
\left[  x^{j}\right]  f=\left[  x^{j}\right]  g
\label{pf.lem.fps.prod.irlv.fg.1}%
\end{equation}
(by (\ref{eq.lem.fps.prod.irlv.fg.ass}), applied to $j$ instead of $m$).
However, (\ref{pf.thm.fps.ring.xn(ab)=3}) (applied to $m$, $a$ and $f$ instead
of $n$, $\mathbf{a}$ and $\mathbf{b}$) yields%
\[
\left[  x^{m}\right]  \left(  af\right)  =\sum_{j=0}^{m}\left[  x^{m-j}%
\right]  a\cdot\underbrace{\left[  x^{j}\right]  f}_{\substack{=\left[
x^{j}\right]  g\\\text{(by (\ref{pf.lem.fps.prod.irlv.fg.1}))}}}=\sum
_{j=0}^{m}\left[  x^{m-j}\right]  a\cdot\left[  x^{j}\right]  g.
\]
On the other hand, (\ref{pf.thm.fps.ring.xn(ab)=3}) (applied to $m$, $a$ and
$g$ instead of $n$, $\mathbf{a}$ and $\mathbf{b}$) yields%
\[
\left[  x^{m}\right]  \left(  ag\right)  =\sum_{j=0}^{m}\left[  x^{m-j}%
\right]  a\cdot\left[  x^{j}\right]  g.
\]
Comparing these two equalities, we obtain $\left[  x^{m}\right]  \left(
af\right)  =\left[  x^{m}\right]  \left(  ag\right)  $. This proves Lemma
\ref{lem.fps.prod.irlv.fg}.
\end{proof}

A consequence of Lemma \ref{lem.fps.prod.irlv.fg} is the following fact:

\begin{lemma}
\label{lem.fps.prod.irlv.mul}Let $u,v\in K\left[  \left[  x\right]  \right]  $
be two FPSs such that $v$ is a multiple of $u$. Let $n\in\mathbb{N}$. Assume
that
\begin{equation}
\left[  x^{m}\right]  u=0\ \ \ \ \ \ \ \ \ \ \text{for each }m\in\left\{
0,1,\ldots,n\right\}  . \label{eq.lem.fps.prod.irlv.mul.ass}%
\end{equation}
Then,
\[
\left[  x^{m}\right]  v=0\ \ \ \ \ \ \ \ \ \ \text{for each }m\in\left\{
0,1,\ldots,n\right\}  .
\]

\end{lemma}

\begin{proof}
[Proof of Lemma \ref{lem.fps.prod.irlv.mul}.]We have assumed that $v$ is a
multiple of $u$. In other words, $v=ua$ for some $a\in K\left[  \left[
x\right]  \right]  $. Consider this $a$.

For each $m\in\left\{  0,1,\ldots,n\right\}  $, we have%
\begin{align*}
\left[  x^{m}\right]  u  &  =0\ \ \ \ \ \ \ \ \ \ \left(  \text{by
(\ref{eq.lem.fps.prod.irlv.mul.ass})}\right) \\
&  =\left[  x^{m}\right]  0\ \ \ \ \ \ \ \ \ \ \left(  \text{since the FPS
}0\text{ satisfies }\left[  x^{m}\right]  0=0\right)  .
\end{align*}
Hence, Lemma \ref{lem.fps.prod.irlv.fg} (applied to $f=u$ and $g=0$) yields
that%
\begin{equation}
\left[  x^{m}\right]  \left(  au\right)  =\left[  x^{m}\right]  \left(
a\cdot0\right)  \ \ \ \ \ \ \ \ \ \ \text{for each }m\in\left\{
0,1,\ldots,n\right\}  . \label{pf.lem.fps.prod.irlv.mul.2}%
\end{equation}
Now, for each $m\in\left\{  0,1,\ldots,n\right\}  $, we have%
\begin{align*}
\left[  x^{m}\right]  v  &  =\left[  x^{m}\right]  \left(  au\right)
\ \ \ \ \ \ \ \ \ \ \left(  \text{since }v=ua=au\right) \\
&  =\left[  x^{m}\right]  \underbrace{\left(  a\cdot0\right)  }_{=0}%
\ \ \ \ \ \ \ \ \ \ \left(  \text{by (\ref{pf.lem.fps.prod.irlv.mul.2}%
)}\right) \\
&  =\left[  x^{m}\right]  0=0.
\end{align*}
This proves Lemma \ref{lem.fps.prod.irlv.mul}.
\end{proof}

We can derive a further useful consequence from Lemma
\ref{lem.fps.prod.irlv.mul}:

\begin{lemma}
\label{lem.fps.prod.irlv.cong-mul}Let $a,b,c,d\in K\left[  \left[  x\right]
\right]  $ be four FPSs. Let $n\in\mathbb{N}$. Assume that
\begin{equation}
\left[  x^{m}\right]  a=\left[  x^{m}\right]  b\ \ \ \ \ \ \ \ \ \ \text{for
each }m\in\left\{  0,1,\ldots,n\right\}  .
\label{pf.lem.fps.prod.irlv.cong-mul.assab}%
\end{equation}
Assume further that%
\begin{equation}
\left[  x^{m}\right]  c=\left[  x^{m}\right]  d\ \ \ \ \ \ \ \ \ \ \text{for
each }m\in\left\{  0,1,\ldots,n\right\}  .
\label{pf.lem.fps.prod.irlv.cong-mul.asscd}%
\end{equation}
Then,
\[
\left[  x^{m}\right]  \left(  ac\right)  =\left[  x^{m}\right]  \left(
bd\right)  \ \ \ \ \ \ \ \ \ \ \text{for each }m\in\left\{  0,1,\ldots
,n\right\}  .
\]

\end{lemma}

\begin{proof}
[Proof of Lemma \ref{lem.fps.prod.irlv.cong-mul}.]For each $m\in\left\{
0,1,\ldots,n\right\}  $, we have%
\begin{align*}
\left[  x^{m}\right]  \left(  a-b\right)   &  =\left[  x^{m}\right]  a-\left[
x^{m}\right]  b\ \ \ \ \ \ \ \ \ \ \left(  \text{by
(\ref{pf.thm.fps.ring.xn(a-b)=})}\right) \\
&  =0\ \ \ \ \ \ \ \ \ \ \left(  \text{by
(\ref{pf.lem.fps.prod.irlv.cong-mul.assab})}\right)  .
\end{align*}
Moreover, the FPS $ac-bc$ is a multiple of $a-b$ (since $ac-bc=\left(
a-b\right)  c$). Hence, Lemma \ref{lem.fps.prod.irlv.mul} (applied to $u=a-b$
and $v=ac-bc$) shows that
\begin{equation}
\left[  x^{m}\right]  \left(  ac-bc\right)  =0\ \ \ \ \ \ \ \ \ \ \text{for
each }m\in\left\{  0,1,\ldots,n\right\}
\label{pf.lem.fps.prod.irlv.cong-mul.2}%
\end{equation}
(since we have $\left[  x^{m}\right]  \left(  a-b\right)  =0$ for each
$m\in\left\{  0,1,\ldots,n\right\}  $).

For each $m\in\left\{  0,1,\ldots,n\right\}  $, we have%
\begin{align*}
\left[  x^{m}\right]  \left(  c-d\right)   &  =\left[  x^{m}\right]  c-\left[
x^{m}\right]  d\ \ \ \ \ \ \ \ \ \ \left(  \text{by
(\ref{pf.thm.fps.ring.xn(a-b)=})}\right) \\
&  =0\ \ \ \ \ \ \ \ \ \ \left(  \text{by
(\ref{pf.lem.fps.prod.irlv.cong-mul.asscd})}\right)  .
\end{align*}
Moreover, the FPS $bc-bd$ is a multiple of $c-d$ (since $bc-bd=b\left(
c-d\right)  =\left(  c-d\right)  b$). Hence, Lemma \ref{lem.fps.prod.irlv.mul}
(applied to $u=c-d$ and $v=bc-bd$) shows that
\begin{equation}
\left[  x^{m}\right]  \left(  bc-bd\right)  =0\ \ \ \ \ \ \ \ \ \ \text{for
each }m\in\left\{  0,1,\ldots,n\right\}
\label{pf.lem.fps.prod.irlv.cong-mul.4}%
\end{equation}
(since we have $\left[  x^{m}\right]  \left(  c-d\right)  =0$ for each
$m\in\left\{  0,1,\ldots,n\right\}  $).

Now, let $m\in\left\{  0,1,\ldots,n\right\}  $. Then,
(\ref{pf.thm.fps.ring.xn(a-b)=}) yields $\left[  x^{m}\right]  \left(
ac-bc\right)  =\left[  x^{m}\right]  \left(  ac\right)  -\left[  x^{m}\right]
\left(  bc\right)  $. Comparing this with
(\ref{pf.lem.fps.prod.irlv.cong-mul.2}), we obtain $\left[  x^{m}\right]
\left(  ac\right)  -\left[  x^{m}\right]  \left(  bc\right)  =0$. In other
words, $\left[  x^{m}\right]  \left(  ac\right)  =\left[  x^{m}\right]
\left(  bc\right)  $. On the other hand, (\ref{pf.thm.fps.ring.xn(a-b)=})
yields $\left[  x^{m}\right]  \left(  bc-bd\right)  =\left[  x^{m}\right]
\left(  bc\right)  -\left[  x^{m}\right]  \left(  bd\right)  $. Comparing this
with (\ref{pf.lem.fps.prod.irlv.cong-mul.4}), we obtain $\left[  x^{m}\right]
\left(  bc\right)  -\left[  x^{m}\right]  \left(  bd\right)  =0$. In other
words, $\left[  x^{m}\right]  \left(  bc\right)  =\left[  x^{m}\right]
\left(  bd\right)  $. Hence, $\left[  x^{m}\right]  \left(  ac\right)
=\left[  x^{m}\right]  \left(  bc\right)  =\left[  x^{m}\right]  \left(
bd\right)  $. This proves Lemma \ref{lem.fps.prod.irlv.cong-mul}.
\end{proof}
